% Companion to kg_unified_multimodal.tex: how the embedding matrix E and relation
% parameters Theta are allocated, initialized, and fitted (not assumed given).
% Build: pdflatex kg_embedding_construction.tex
\documentclass[11pt]{article}

\usepackage{amsmath,amssymb,amsthm,mathtools}
\usepackage{bm}
\usepackage[margin=1in]{geometry}
\usepackage{hyperref}
\hypersetup{colorlinks=true,linkcolor=blue,citecolor=blue,urlcolor=blue}

\newtheorem{definition}{Definition}
\newtheorem{remark}{Remark}
\newtheorem{proposition}{Proposition}

\title{Constructing the Geometric Layer:\\
Initialization, Parameter Blocks, and Optimization for Knowledge Graph Embeddings}
\author{Niklaus Parcell\\
\textit{Mont Ops Inc. / Independent Research}\\
\texttt{nik@mont-ops.com}}
\date{\today}

\begin{document}
\maketitle

\begin{abstract}
The linear-algebraic framework of Parcell \cite{UnifiedNote} treats entity rows $\bm{e}_a$ and relation parameters $\bm{\theta}_r$ as inputs to a scoring map $s(h,r,t)$.
In deployment, those dense parameters are rarely provided \emph{ab initio}; they must be \emph{constructed} by allocation, optional structured initialization, and (most often) iterative optimization on link-prediction or ranking objectives coupled with negative sampling.
This companion note formalizes the \textbf{parameter bundle} $(\bm{E}, \bm{\Theta})$, standard initialization choices, a generic stochastic training loop, and complexity at the level of scoring forward passes---complementing the structural (sparse) layer developed in the unified note.
\end{abstract}

\section{Relationship to the unified framework}
The unified note (henceforth \textsc{Unified}) separates a sparse structural layer derived from asserted triples $\mathcal{T}$ from a geometric layer of embeddings and scores.
\textsc{Unified} is largely agnostic to \emph{how} $\bm{E} \in \mathbb{R}^{N \times d}$ and relation parameters $\bm{\Theta}$ arise.
This note closes that gap: $\bm{E}$ and $\bm{\Theta}$ are \emph{outputs} of a construction pipeline (allocate $\to$ initialize $\to$ optionally train $\to$ persist), not prerequisites handed in by an oracle.

\section{Structured parameter bundle}

\subsection{Entity embedding table}
\begin{definition}[Entity table]
The \emph{entity embedding matrix} is $\bm{E} \in \mathbb{R}^{N \times d}$, whose $a$-th row $\bm{e}_a^{\mathsf{T}}$ is the embedding of entity $a \in [N]$.
Storage is typically \emph{row-major} dense layout in numerical libraries (each row contiguous).
\end{definition}

\subsection{Relation parameter blocks}
\begin{definition}[Relation parameterization]
A \emph{relation block} for type $r \in [L]$ is an ordered collection of tensors $\bm{\theta}_r$ (vectors, matrices, or higher-order blocks) whose shapes are fixed by the scoring architecture.
The full \emph{relation bundle} is $\bm{\Theta} = (\bm{\theta}_1,\ldots,\bm{\theta}_L)$.
\end{definition}

\paragraph{Custom structure.}
Implementation-wise, $\bm{\Theta}$ need not flatten into a single matrix; it may be a struct of arrays (e.g., one dense $\bm{W}_r \in \mathbb{R}^{d \times d}$ per relation for DistMult-style models, or lower-rank factors).
What matters algorithmically is a deterministic map
\[
  (\bm{E}, \bm{\Theta}) \;\mapsto\; s(h,r,t \,;\, \bm{E}, \bm{\Theta})
\]
differentiable or subdifferentiable wherever optimization demands.

\begin{remark}[Alignment with \textsc{Unified}]
\textsc{Unified} writes scores abstractly; this note treats $(\bm{E}, \bm{\Theta})$ as the mutable state vectorized logically (if not physically) for optimizers.
Sparse operators $\bm{L}_k$ from $\mathcal{T}$ act on \emph{current} node features (often initial rows of $\bm{E}$ or derived $\bm{H}$); training may alternate graph propagation and score-based losses.
\end{remark}

\section{Allocation and initialization}

\subsection{Allocation}
Given $(N, L, d)$ and the per-relation shapes implied by the model class, an implementation:
\begin{enumerate}
  \item allocates $\bm{E}$ with $Nd$ scalars;
  \item allocates each $\bm{\theta}_r$ per schema;
  \item optionally reserves optimizer state (momentum, second moments) of matching shapes.
\end{enumerate}

\subsection{Initialization strategies}
Common choices (not mutually exclusive across ablations):
\begin{itemize}
  \item \textbf{Random isotropic / fan schemes}: entries i.i.d.\ or scaled by layer width to control gradient magnitude at start.
  \item \textbf{Small-norm}: enforce $\lVert \bm{e}_a \rVert$ bounds initially for translation models.
  \item \textbf{Deterministic prototype}: embed entity $a$ from hashed metadata or degree statistics (weak prior, debugging).
  \item \textbf{Warm start}: load $(\bm{E}, \bm{\Theta})$ from checkpoint or external encoder, then continue training.
  \item \textbf{Non-learned placeholder}: $\bm{E}$ left fixed (e.g., zeros or one-hot padded); useful only for structural software tests.
\end{itemize}

\begin{remark}
No initialization implies factual correctness; only optimization (or external grounding) aligns $\bm{E}$ with $\mathcal{T}$ and downstream tasks.
\end{remark}

\section{Negative sampling and training batches}

\subsection{Positive minibatch}
Fix batch size $B$.
A minibatch $\mathcal{B}^+ = \{ (h_b, r_b, t_b) \}_{b=1}^{B}$ samples asserted triples (with replacement or without, from empirical distribution over $\mathcal{T}$).

\subsection{Corrupted negatives}
For each positive $(h,r,t)$, generate one or more negatives by replacing $t$ (or $h$) with an index drawn from a distribution $q$ over $[N]$ (uniform, frequency-based, or adversarial).
Denote negatives $(h', r', t')$.

\subsection{Objectives (sketch)}
\paragraph{Margin ranking.}
\begin{equation}
  \mathcal{L}_{\mathrm{margin}} = \sum_{b} \max\bigl(0,\, \gamma - s(h_b,r_b,t_b) + s(h_b^{\text{neg}}, r_b, t_b^{\text{neg}}) \bigr).
\end{equation}

\paragraph{Binary logistic / softplus.}
Scores against sigmoid of $s_{\mathrm{pos}} - s_{\mathrm{neg}}$.

\paragraph{Multi-class / sampled softmax.}
Treat tail prediction among a candidate set including the true tail and $K$ negatives.

\section{Generic construction algorithm}

\begin{definition}[Construction pipeline]
A \emph{construction pipeline} is a tuple $(\mathrm{Init}, \mathrm{Train?}, \mathrm{Save})$ where $\mathrm{Init}$ produces $(\bm{E}^{(0)}, \bm{\Theta}^{(0)})$, $\mathrm{Train?}$ optionally maps $T$ steps to $(\bm{E}^{(T)}, \bm{\Theta}^{(T)})$, and $\mathrm{Save}$ serializes tensors.
\end{definition}

\paragraph{Training loop (conceptual).}
For step $t = 0,\ldots,T-1$:
\begin{enumerate}
  \item Sample minibatch $\mathcal{B}^+$ from $\mathcal{T}$; build negatives.
  \item Evaluate forward scores $s$ for positives and negatives using $(\bm{E}^{(t)}, \bm{\Theta}^{(t)})$.
  \item Compute loss $\mathcal{L}^{(t)}$ (optionally + regularizers on $\lVert \bm{E} \rVert_F$ or relation blocks).
  \item Update parameters via a chosen optimizer (SGD, Adam, \ldots), respecting numeric dtype \texttt{f32}/\texttt{f64}.
\end{enumerate}

\begin{proposition}[Per-step dominant cost (forward)]
\label{prop:batch-cost}
Suppose scoring each triple requires $\Theta(d)$ or $\Theta(d^2)$ work for bilinear models with $d \times d$ relation matrices, and batch size is $B$ with $K$ negatives per positive.
Then one minibatch forward through scores is $\Theta(B(1+K)\cdot C(d))$ where $C(d)$ is the per-triple asymptotic score cost.
Sparse propagations $\bm{L}\bm{X}$ incur additional $O(Zd)$ cost per application as in \textsc{Unified}.
\end{proposition}

\begin{proof}[Proof sketch]
Count scalar operations in $s$ for one triple under the chosen model; multiply by the number of scored triples in the batch.
\end{proof}

\section{Persistence and versioning}
Serialize $(N,L,d)$, vocabulary maps, $\bm{E}$, flattened or structured $\bm{\Theta}$, optimizer metadata, and training configuration.
Checksums or semantic versioning protect consumer code from silent shape mismatch.

\section{Discussion}
This note deliberately omits dataset-specific hyperparameter recipes and autodifferentiation plumbing; it specifies \emph{what} must exist for $\bm{E}$ to be ``created'' in software: allocation, initialization, optional iterative fitting driven by $(\mathcal{T}, s, \mathcal{L})$, and persistence.
Multimodal alignment losses add terms coupling rows of $\bm{E}$ to projected encoder outputs; structurally they extend $\mathcal{L}$ without changing the allocation story.

\begin{thebibliography}{9}

\bibitem{UnifiedNote}
N.~Parcell.
\newblock \emph{Sparse Knowledge Graphs, Joint Embeddings, and Multimodal Alignment: A Linear-Algebraic Framework for Structured World Models.}
\newblock Companion technical note \texttt{kg\_unified\_multimodal.tex}, 2026.

\bibitem{Bordes2013}
A.~Bordes, N.~Usunier, A.~Garc\'ia-Dur\'an, J.~Weston, O.~Yakhnenko.
\newblock Translating embeddings for modeling multi-relational data.
\newblock In \emph{NIPS}, 2013.

\bibitem{Niu2025}
G.~Niu et al.
\newblock Knowledge graph embeddings: a comprehensive survey on capturing relation properties.
\newblock \emph{arXiv preprint} 2410.14733, 2025.

\end{thebibliography}

\end{document}
